%Version 3.1 December 2024
\documentclass[pdflatex,sn-mathphys-num]{sn-jnl}

\usepackage{graphicx}
\usepackage{float}
\usepackage[section]{placeins}
\usepackage{amsmath,amssymb,amsfonts}
\usepackage{amsthm}
\usepackage{mathrsfs}
\usepackage[title]{appendix}
\usepackage{xcolor}
\usepackage{textcomp}
\usepackage{manyfoot}
\usepackage{booktabs}
\usepackage{algorithm}
\usepackage{algorithmicx}
\usepackage{algpseudocode}
\usepackage{listings}
\usepackage{nicefrac}
\usepackage{xspace}

\DeclareRobustCommand{\mathdefault}[1]{#1}

\newcommand{\vMLabel}{[\textbf{vM}]\xspace}
\newcommand{\Kinc}{k_{\rm inc}}
\newcommand{\Gcinc}{G_{c,\rm inc}}
\newcommand{\GcZero}{G_c^0}
\newcommand{\Jeff}{J_{x,c}^{\rm eff}}

\newcommand{\NumCases}{45}
\newcommand{\WstegValues}{0.4, 0.5, 0.6, 1, 2}
\newcommand{\JeffBaselineMin}{1.75}
\newcommand{\JeffBaselineMax}{1.80}
\newcommand{\MeanGcLowEffect}{-27.5\%}
\newcommand{\MeanGcHighEffect}{-10.0\%}
\newcommand{\MeanKLowEffect}{-25.8\%}
\newcommand{\MeanKHighEffect}{-14.3\%}


\raggedbottom

\begin{document}

\title[Effective crack resistance of a ductile material with circular inclusions]{Effective crack resistance of a simplified ductile heterogeneous material with circular inclusions}

\author*[1]{\fnm{Alexander} \sur{Schl\"{u}ter}}
\email{alexander.schlueter@tu-darmstadt.de}

\author[1]{\fnm{Ronjit} \sur{Medda}}
\author[1]{\fnm{Ralf} \sur{M\"{u}ller}}

\affil*[1]{%
  \orgdiv{Institut f\"ur Mechanik, FB Bau- und Umweltingenieurwissenschaften},
  \orgname{TU Darmstadt},
  \orgaddress{%
    \street{Franziska-Braun-Stra{\ss}e 7},
    \city{Darmstadt},
    \postcode{64287},
    \country{Germany}
  }
}

\abstract{
This short contribution investigates the effective crack resistance obtained from a numerical crack-growth experiment for a simplified two-dimensional ductile heterogeneous material. In contrast to previous studies on pores, the resolved microstructure contains circular inclusions. Their stiffness and fracture resistance are varied independently together with the ligament width between neighboring inclusions. Crack evolution is simulated with a phase-field model for ductile fracture based on von Mises plasticity. The maximum boundary J-integral during enforced macroscopic crack advance is used as the operational measure of effective crack resistance. The results show that inclusion toughness and inclusion stiffness are strongly coupled: weak inclusions reduce the effective crack resistance, whereas tough inclusions do not produce a universal increase. The largest responses in the retained parameter range occur for tough compliant inclusions and for stiff weak inclusions, showing that the crack path is as important as the local fracture-resistance value.
}

\keywords{Ductile fracture, Phase-field fracture, Heterogeneous materials, Circular inclusions, Effective crack resistance, J-integral}

\maketitle

\section{Introduction}\label{sec:intro}

The determination of fracture resistance parameters for heterogeneous materials remains less mature than the homogenization of elastic and plastic stiffness. Classical fracture mechanics introduces critical quantities such as the stress-intensity factor \(K_{Ic}\), the Griffith energy release rate \(G_c\), or, more generally, a critical value \(J_c\) in a criterion
\begin{equation}
    J = J_c .
    \label{eq:fracture_criterion_J}
\end{equation}
For homogeneous ductile materials this type of energetic criterion remains meaningful at least up to the onset of monotonic crack growth, cf. Gross and Seelig~\cite{gross_fracture_2011}. For heterogeneous materials, however, it is not obvious how a single macroscopic fracture parameter should be defined, whether it converges under enlargement of a representative domain, and how it depends on the resolved microstructure.

This question is particularly relevant for additively manufactured metallic materials, where pores, inclusions, and locally different material regions may be introduced by the process or by an intentional multi-material design. Wire arc additive manufacturing, for instance, enables the fabrication of complex metallic structures and also motivates fracture-mechanics-based assessment in materials with non-trivial microstructure~\cite{hauser_porosity_2021}. A practical macroscopic analysis then requires parameters that condense the fracture process at the microscale into a usable crack-resistance measure.

Different routes toward effective fracture properties have been proposed. Schneider~\cite{schneider_fftbased_2020} and Michel and Suquet~\cite{michel_merits_2022}, among others, discuss variational definitions related to the surface energy of a heterogeneous medium. In the present work we follow the numerical experiment introduced by Hossain et al.~\cite{hossain_effective_2014}. A representative material region is forced to separate under surfing boundary conditions, while the macroscopic energetic driving force is measured by a boundary J-integral. The maximum value attained during the process is interpreted as an effective crack resistance for this particular experiment. This terminology is used deliberately: the value is an operational crack-resistance measure for the prescribed microstructure and loading path, not a proof of a homogenized fracture toughness.

Previous work using this methodology focused on porous materials and compared different elastic and plastic constitutive descriptions. Here the aim is narrower and more direct. We consider a simplified two-dimensional material with circular inclusions instead of pores. The matrix follows a von Mises-type ductile phase-field fracture model, while the inclusions are assigned independently varied stiffness and fracture resistance. The study asks how these inclusion properties and the ligament width between inclusions affect the computed effective crack resistance. The manuscript is kept short: Section~\ref{sec:numericalExperiment} recalls the numerical experiment, Section~\ref{sec:simulation_phase_field_models} summarizes the ductile phase-field model and implementation, Section~\ref{sec:2Dresults} presents the inclusion study, and Section~\ref{sec:discussion} discusses the observed trends.

\section{A Numerical Experiment to Determine Effective Crack Resistance}\label{sec:numericalExperiment}

\begin{figure}[H]
    \centering
    \resizebox{0.6\textwidth}{!}{\input{pics/num_experiment.pdf_t}}
    \caption{Sketch of the numerical experiment introduced by Hossain et al.~\cite{hossain_effective_2014}. A material region with explicitly resolved microstructure is separated in a simulation while the boundary J-integral records the macroscopic crack-driving force.}
    \label{fig:numerical_experiment}
\end{figure}

The numerical experiment is shown schematically in Fig.~\ref{fig:numerical_experiment}. A region containing the resolved microstructure is embedded in a surrounding effective medium. Macroscopic mode-I crack growth is enforced through surfing boundary conditions, where a fictitious crack tip travels through the domain with velocity \(\dot{x}_{\rm bc}\). The boundary displacement field is prescribed from the asymptotic mode-I crack-tip solution,
\begin{align}
    u_x &= \frac{K_I^{\rm bc}}{2\mu_{\rm bc}} \sqrt{\frac{r}{2\pi}}\left(\kappa-\cos\varphi\right)\cos\frac{\varphi}{2},\\
    u_y &= \frac{K_I^{\rm bc}}{2\mu_{\rm bc}} \sqrt{\frac{r}{2\pi}}\left(\kappa-\cos\varphi\right)\sin\frac{\varphi}{2},
    \label{eq:surfing-bcs}
\end{align}
with \(\kappa=(3-\nu_{\rm bc})/(1+\nu_{\rm bc})\) and \(\mu_{\rm bc}=E_{\rm bc}/(2(1+\nu_{\rm bc}))\). Here \(r\) and \(\varphi\) are polar coordinates centered at the fictitious crack tip. The microscopic crack path is not prescribed; it follows from the phase-field fracture simulation.

During the separation process, the macroscopic driving force is evaluated as the \(x\)-component of the J-integral at the outer boundary,
\begin{equation}
    J_x = \int_{\partial\Omega}\left(\psi_{\rm el}\boldsymbol{1} - (\nabla\boldsymbol{u})^T\boldsymbol{\sigma}\right)\boldsymbol{n}\,{\rm d}A \cdot \boldsymbol{e}_x .
    \label{eq:J-integral}
\end{equation}
The effective crack resistance used in this work is then
\begin{equation}
    \Jeff = \max_t J_x(t).
    \label{eq:effective_resistance}
\end{equation}
This definition captures the largest macroscopic energy flux required to force complete separation of the resolved material region. It is therefore a property of the chosen numerical experiment, microstructure, constitutive model, and crack-growth direction.

\section{Simulation of Fracture with Phase Field Models}\label{sec:simulation_phase_field_models}

\subsection{Ductile phase-field fracture model}\label{sec:phase_fiel_model_ductile}

Fracture in the ductile matrix is modeled by the phase-field formulation of Noll et al.~\cite{noll_3d_2020}. Small strains are assumed and the total strain is additively decomposed as
\begin{equation}
    \boldsymbol{\varepsilon} = \boldsymbol{\varepsilon}^e+\boldsymbol{\varepsilon}^p .
\end{equation}
The unfractured material follows von Mises plasticity with linear isotropic hardening. The yield stress is
\begin{equation}
    \sigma_Y(\alpha)=\sigma_y+H\alpha,
\end{equation}
where \(\alpha\) is the accumulated plastic variable, \(H\) the hardening modulus, and \(\sigma_y\) the initial yield stress. With \(\boldsymbol{s}=\operatorname{dev}(\boldsymbol{\sigma})\), the yield function reads
\begin{equation}
    f(\boldsymbol{\sigma},\alpha)=\|\boldsymbol{s}\|-\sqrt{\frac{2}{3}}\left(\sigma_y+H\alpha\right).
\end{equation}
The evolution of the plastic variables follows the usual associative \(J_2\) flow rule together with the Karush--Kuhn--Tucker conditions.

Fracture is described by a scalar phase field \(s(\boldsymbol{x},t)\in[0,1]\), where \(s=1\) denotes intact material and \(s=0\) fully fractured material. The degraded elastic energy density is
\begin{equation}
    \psi_{\rm el}(\boldsymbol{\varepsilon},\boldsymbol{\varepsilon}^p,s)
    = g(s)\frac{1}{2}
    \left(\boldsymbol{\varepsilon}-\boldsymbol{\varepsilon}^p\right)
    :
    \left(\boldsymbol{\mathbb{C}}:
    \left(\boldsymbol{\varepsilon}-\boldsymbol{\varepsilon}^p\right)\right),
\end{equation}
with the cubic degradation function
\begin{equation}
    g(s)=\beta(s^3-s^2)+3s^2-2s^3+\eta .
    \label{eq:deg_grad}
\end{equation}
In this work \(\beta=0.2\) and \(\eta=10^{-5}\). The plastic contribution is degraded in the same manner,
\begin{equation}
    \psi_{\rm pl}(\alpha,s)=g(s)\left(\frac{1}{2}H\alpha^2+\sigma_y\alpha\right).
\end{equation}
The regularized fracture energy density is
\begin{equation}
    \psi_{\rm fr}=G_c\left(\frac{(1-s)^2}{2\epsilon}+2\epsilon\|\nabla s\|^2\right),
\end{equation}
where \(G_c\) is the local fracture resistance and \(\epsilon\) controls the width of the diffuse crack zone. The resulting phase-field evolution equation is
\begin{align}
    \dot{s}=-M\left[
    g'(s)\left(\psi_{\rm el}+\left(\sigma_y+\frac{1}{2}H\alpha\right)\alpha\right)
    -G_c\left(2\epsilon\Delta s+\frac{1-s}{2\epsilon}\right)
    \right],
    \label{eq:evolution_s}
\end{align}
with homogeneous Neumann boundary conditions for \(s\). Quasi-static mechanical equilibrium is imposed by
\begin{equation}
    \operatorname{div}\boldsymbol{\sigma}=\boldsymbol{0}.
    \label{eq:equilibrium}
\end{equation}
The stress is computed as
\begin{equation}
    \boldsymbol{\sigma}=g(s)\left[
    K\,\operatorname{tr}(\boldsymbol{\varepsilon})\boldsymbol{1}
    +2\mu\left(\operatorname{dev}(\boldsymbol{\varepsilon})-\boldsymbol{\varepsilon}^p\right)
    \right].
    \label{eq:sigma_plast}
\end{equation}
The matrix material parameters used throughout the inclusion study are \(\lambda=\mu=1.0\), \(G_c=\GcZero=1.0\), \(\sigma_y=1.0\), and \(H=0.2222222\). The inclusion yield stress is set to \(100\sigma_y\), so that the inclusions remain effectively elastic in the present simulations.

\subsection{Numerical implementation}\label{sec:numerical_implementation}

The weak forms are implemented in FEniCSx~\cite{_fenicsx_}. The displacement, phase field, and plastic internal variables are solved monolithically in pseudo-time. The phase-field equation is discretized by a backward Euler scheme, and the nonlinear system is solved by Newton's method. If convergence is not reached within the prescribed iteration limit, the time step is reduced and the load step is repeated.

First-order Lagrange finite elements are used on triangular meshes. The mesh is refined in the expected fracture region such that the phase-field length scale is resolved. In the representative stored simulation used for visualization, the minimum edge length is \(1.30\times 10^{-2}L\), the mean edge length is \(3.24\times 10^{-2}L\), and \(\epsilon=0.1L\). The graph data used for the present evaluation were written by the simulation code as scalar histories containing pseudo-time, \(J_x\), \(J_y\), crack-tip position \(x_{\rm ct}\), and the surfing-boundary-condition position \(x_{\rm bc}\).

\section{Effective Crack Resistance of a Simplified Two-Dimensional Heterogeneous Material}\label{sec:2Dresults}

\subsection{Geometry and material parameters}\label{sec:geometry}

\begin{figure}[H]
    \centering
    \begin{minipage}{0.60\textwidth}
        \centering
        \textbf{a)}\\[-0.3em]
        \makebox[\textwidth][c]{\includegraphics[width=\textwidth]{pics/material_layout_inclusions.png}}
    \end{minipage}
    \par\medskip
    \begin{minipage}{0.60\textwidth}
        \centering
        \textbf{b)}\\[-0.3em]
        \makebox[\textwidth][c]{\includegraphics[width=\textwidth]{pics/failure_pattern_phasefield.png}}
    \end{minipage}
    \par\medskip
    \begin{minipage}{0.60\textwidth}
        \centering
        \textbf{c)}\\[-0.3em]
        \makebox[\textwidth][c]{\includegraphics[width=\textwidth]{pics/mesh_representative.png}}
    \end{minipage}
    \caption{Representative setup for the inclusion study. a) Full computational domain, including the resolved material region with circular inclusions and the surrounding effective medium. b) Final phase-field pattern for one stored simulation, plotted with a blue-to-red scale for the phase field \(s\). c) Corresponding triangular finite-element mesh. The shown simulation has \(w_s=1.0L\), \(\Kinc=0.5\), and \(\Gcinc=0.5\GcZero\).}
    \label{fig:setup_inclusions}
\end{figure}

The investigated two-dimensional material contains circular inclusions of diameter \(L\) placed in a regular array. The distance between neighboring inclusion boundaries is denoted by \(w_s\), so that the cell size is \(w_c=L+w_s\). The resolved region contains three rows and six columns of inclusions and is embedded in a surrounding effective medium. Plane strain conditions are assumed.

The parameter study varies the ligament width
\[
    w_s/L \in \{\WstegValues\},
\]
the inclusion stiffness factor
\[
    \Kinc \in \{0.5,1.0,1.5\},
\]
and the inclusion fracture-resistance factor
\[
    \Gcinc/\GcZero \in \{0.5,1.0,1.5\}.
\]
The inclusion Lam\'e parameters are obtained by scaling the matrix values as
\begin{equation}
    \lambda_{\rm inc}=\Kinc\lambda_{\rm mat}, \qquad
    \mu_{\rm inc}=\Kinc\mu_{\rm mat},
\end{equation}
and the inclusion fracture resistance is
\begin{equation}
    G_{c,\rm inc} = (\Gcinc/\GcZero)\,\GcZero .
\end{equation}
For each geometry and stiffness contrast, a preliminary effective-stiffness computation provides the Lam\'e parameters of the surrounding homogeneous material. A total of \(\NumCases\) complete readable simulation histories were available in the current result folder and were used for the plots below. The cases with \(w_s=0.75L\) were excluded from the evaluation because this data point showed problematic behavior in the local result snapshot; datasets with \(w_s>2.0L\) were also omitted to focus on the interaction-dominated ligament-width range.

\subsection{Crack tracking and local J-integral response}

\begin{figure}[H]
    \centering
    \includegraphics[width=0.55\textwidth]{pics/fig09_crack_tip_tracking.png}
    \caption{Position of the surfing boundary condition \(x_{\rm bc}\) and extracted crack-tip position \(x_{\rm ct}\) for the reference inclusion case \(w_s=1.0L\), \(\Kinc=1.0\), and \(\Gcinc=\GcZero\).}
    \label{fig:crack_tip_tracking}
\end{figure}

Figure~\ref{fig:crack_tip_tracking} shows that the measured crack tip follows the moving boundary-condition field after the initial crack-growth transient. Small plateaus and jumps in \(x_{\rm ct}\) are expected because the crack interacts with inclusions and because the phase-field crack tip is extracted from a diffuse crack representation.

\begin{figure}[H]
    \centering
    \includegraphics[width=0.72\textwidth]{pics/fig10_Jx_xct_constant_Gc_vary_Kinc.png}
    \caption{J-integral over crack-tip position for \(w_s=1.0L\) and constant inclusion fracture resistance \(\Gcinc=\GcZero\). The inclusion stiffness is varied from compliant to stiff.}
    \label{fig:Jx_xct_stiffness}
\end{figure}

For fixed \(w_s=1.0L\) and \(\Gcinc=\GcZero\), Fig.~\ref{fig:Jx_xct_stiffness} isolates the effect of inclusion stiffness over the full recorded crack-growth history. The stiffness contrast changes the local peak sequence and the crack-advance plateaus, but the effect is not a simple monotonic shift. This indicates that stiffness contrast changes more than the local elastic field around the inclusion: it also modifies how the crack advances through the heterogeneous region.

\begin{figure}[H]
    \centering
    \includegraphics[width=0.9\textwidth]{pics/fig11_Jx_sections_Gcinc_0p5.png}
    \caption{Section of the J-integral \(J_x\) versus crack-tip position \(x_{\rm ct}\) for \(\Gcinc=0.5\GcZero\). Only the characteristic evaluation range determined with \(starting\_hole\_to\_evaluate=4\) is shown; for \(w_s=L\), this corresponds approximately to \(6\le x_{\rm ct}/L \le 10\). The three panels show the available simulations for \(\Kinc=0.5\), \(1.0\), and \(1.5\); within each panel the curves correspond to different ligament widths \(w_s\).}
    \label{fig:Jx_sections_Gcinc_0p5}
\end{figure}

\begin{figure}[H]
    \centering
    \includegraphics[width=0.9\textwidth]{pics/fig12_Jx_sections_Gcinc_1p0.png}
    \caption{Section of the J-integral \(J_x\) versus crack-tip position \(x_{\rm ct}\) for \(\Gcinc=\GcZero\). Only the characteristic evaluation range determined with \(starting\_hole\_to\_evaluate=4\) is shown; for \(w_s=L\), this corresponds approximately to \(6\le x_{\rm ct}/L \le 10\). For the stiff-inclusion case \(\Kinc=1.5\), the lower bound of the evaluation interval is shifted by \(w_s\) to exclude the preceding interaction. The three panels show the available simulations for \(\Kinc=0.5\), \(1.0\), and \(1.5\); within each panel the curves correspond to different ligament widths \(w_s\).}
    \label{fig:Jx_sections_Gcinc_1p0}
\end{figure}

\begin{figure}[H]
    \centering
    \includegraphics[width=0.9\textwidth]{pics/fig13_Jx_sections_Gcinc_1p5.png}
    \caption{Section of the J-integral \(J_x\) versus crack-tip position \(x_{\rm ct}\) for \(\Gcinc=1.5\GcZero\). Only the characteristic evaluation range determined with \(starting\_hole\_to\_evaluate=4\) is shown; for \(w_s=L\), this corresponds approximately to \(6\le x_{\rm ct}/L \le 10\). The three panels show the available simulations for \(\Kinc=0.5\), \(1.0\), and \(1.5\); within each panel the curves correspond to different ligament widths \(w_s\). The vertical axis is restricted to \(1.3\le J_x/\GcZero\le2\) to emphasize the relevant peak structure.}
    \label{fig:Jx_sections_Gcinc_1p5}
\end{figure}

Figures~\ref{fig:Jx_sections_Gcinc_0p5}--\ref{fig:Jx_sections_Gcinc_1p5} show only the characteristic J-integral section, rather than the full fracture process. All shown toughness families use \(starting\_hole\_to\_evaluate=4\), consistent with the current evaluation script. The figures make the local crack-inclusion interactions visible: the maxima used in Eq.~\eqref{eq:effective_resistance} are not single smooth material responses, but arise from trapping, deflection, and re-acceleration of the crack as it crosses the inclusion array. The section is taken in the central part of the inclusion row, where the initial transient has decayed. Larger \(w_s\) separates individual interaction events more clearly. The influence of \(\Kinc\) and \(\Gcinc\), however, is not monotonic in all panels, which already indicates that stiffness and local fracture resistance act through the crack path as well as through the local energy scale.

\subsection{Effective crack resistance}

\begin{figure}[H]
    \centering
    \includegraphics[width=0.95\textwidth]{pics/fig15_Jeff_wsteg_all_Gcinc.png}
    \caption{Effective crack resistance \(\Jeff\) versus ligament width \(w_s\). Each panel fixes the inclusion fracture resistance and compares the three inclusion stiffness factors.}
    \label{fig:Jeff_wsteg_families}
\end{figure}

Figure~\ref{fig:Jeff_wsteg_families} summarizes the effective crack resistance for each parameter combination. The ligament width controls the spacing of inclusion-crack interactions and therefore changes the number and magnitude of local resistance peaks. The baseline case \(\Kinc=1.0\), \(\Gcinc=\GcZero\) spans approximately \(\Jeff/\GcZero=\JeffBaselineMin\) to \(\JeffBaselineMax\) over the investigated \(w_s\) range.

At fixed matrix-matched stiffness \(\Kinc=1.0\), changing the inclusion fracture resistance is not symmetric around the baseline. Lowering the inclusion fracture resistance to \(0.5\GcZero\) changes \(\Jeff\) by \(\MeanGcLowEffect\), whereas increasing it to \(1.5\GcZero\) changes \(\Jeff\) by \(\MeanGcHighEffect\). Thus, in the current evaluation window, weak inclusions clearly reduce the effective resistance, but tougher inclusions do not automatically increase the macroscopic value when the stiffness is unchanged.

The inclusion stiffness has an equally path-dependent influence. At \(\Gcinc=\GcZero\), compliant inclusions change the mean effective resistance by \(\MeanKLowEffect\), while stiff inclusions change it by \(\MeanKHighEffect\). Both stiffness contrasts therefore reduce the average value relative to the matrix-matched baseline in the present dataset. In contrast, for weak inclusions \((\Gcinc=0.5\GcZero)\), increasing \(\Kinc\) raises the resistance, while for tough inclusions \((\Gcinc=1.5\GcZero)\), the compliant-inclusion cases carry the largest values. This reversal shows that stiffness contrast cannot be interpreted independently of the inclusion fracture resistance.

\begin{figure}[H]
    \centering
    \includegraphics[width=0.735\textwidth]{pics/fig14_combined_parameter_contribution.png}
    \caption{Combined parameter contribution at selected ligament widths. The entries show \(\Jeff\) normalized by the corresponding baseline case with \(\Kinc=1.0\) and \(\Gcinc=\GcZero\) at the same \(w_s\).}
    \label{fig:combined_parameter_contribution}
\end{figure}

The normalized comparison in Fig.~\ref{fig:combined_parameter_contribution} separates the direct inclusion-property effects from the absolute variation with \(w_s\). The lowest values are obtained for weak and compliant inclusions. The highest values in the present evaluation, however, occur for tough and compliant inclusions rather than for tough and stiff inclusions. The combined effect is therefore not additive: changing stiffness also changes the local crack path and therefore the amount of matrix and inclusion material crossed by the diffuse crack.

\section{Discussion}\label{sec:discussion}

The simulations show that replacing pores by circular inclusions changes the interpretation of the numerical experiment. Pores act primarily as crack traps and stress concentrators, whereas inclusions can either attract, shield, deflect, or resist the crack depending on their stiffness and fracture resistance. In the current evaluation, only \(w_s/L\le2.0\) is retained and the problematic \(w_s=0.75L\) data point is omitted. The reported values therefore focus on the interaction-dominated part of the parameter space rather than on a possible large-ligament saturation regime.

The revised plots show a strong coupling between inclusion toughness and stiffness. Weak inclusions reduce the effective resistance for all stiffnesses, but their response becomes larger as the inclusions are stiffened. Matrix-matched inclusions give the highest resistance for \(\Gcinc=\GcZero\), whereas tough inclusions give the highest resistance when they are compliant. This behavior is consistent with a crack-path effect: the measured maximum J-integral depends not only on the local value of \(G_c\), but also on whether the crack is guided into, around, or between inclusions during the selected central interaction section.

The ligament-width dependence is clearest for the weak-inclusion cases, where \(\Jeff\) generally increases from small \(w_s\) toward \(w_s=2L\). The matrix-matched baseline is comparatively flat, with \(\Jeff/\GcZero\) between \(\JeffBaselineMin\) and \(\JeffBaselineMax\). For tough inclusions the mean response is also nearly flat over the retained ligament widths, but its level depends strongly on stiffness contrast. The rescaled vertical axis in Fig.~\ref{fig:Jx_sections_Gcinc_1p5} makes this especially visible: the relevant peaks remain in a narrow band between \(J_x/\GcZero=1.3\) and \(2\), yet their ordering changes with \(\Kinc\).

Several limitations remain. Only one regular inclusion arrangement, one initial crack position, and one macroscopic crack-growth direction are considered. The result should therefore be read as an inclusion-focused numerical experiment rather than as a general homogenized toughness. In addition, only one representative XDMF/HDF5 field output was available locally for visualization; the quantitative evaluation is based on the scalar history files. Future work should include stochastic inclusion arrangements, larger resolved regions, additional crack directions, and a convergence study with respect to both domain size and phase-field length scale.

\section{Conclusions}\label{sec:conclusions}

A short numerical study was performed to quantify the effective crack resistance of a two-dimensional ductile material with circular inclusions. The phase-field fracture simulations use von Mises plasticity in the matrix and independently varied inclusion stiffness and inclusion fracture resistance. The maximum boundary J-integral during surfing-boundary-condition-driven separation defines the effective crack resistance.

The main findings are:
\begin{itemize}
    \item Weak inclusions reduce the effective crack resistance throughout the retained parameter range.
    \item Tough inclusions do not produce a universal increase; their effect depends strongly on inclusion stiffness and crack path.
    \item Stiffness contrast is coupled to inclusion fracture resistance: stiff inclusions raise the response for weak inclusions, whereas compliant inclusions give the largest response for tough inclusions.
    \item The ligament width remains important because it controls the spacing and intensity of crack-inclusion interactions, especially for weak inclusions.
\end{itemize}
These observations support the use of the numerical experiment as a compact way to compare microstructural design choices for ductile heterogeneous materials.

\bibliography{bib/literatur}

\end{document}
